% Generated draft: requires mathematical and source review.
\documentclass[11pt]{article}
\usepackage[margin=25mm]{geometry}
\usepackage{fontspec}
\setmainfont{Latin Modern Roman}
\usepackage{amsmath,amssymb,mathtools,unicode-math}
\setmathfont{Latin Modern Math}
\usepackage{enumitem,needspace,xurl,hyperref}
\setlength{\parindent}{0pt}
\setlength{\parskip}{5pt}
\setlength{\emergencystretch}{3em}
\newcommand{\N}{\mathbb{N}}
\newcommand{\Z}{\mathbb{Z}}
\newcommand{\Q}{\mathbb{Q}}
\newcommand{\R}{\mathbb{R}}
\newcommand{\C}{\mathbb{C}}
\newcommand{\F}{\mathbb{F}}
\newcommand{\A}{\mathbb{A}}
\newcommand{\PP}{\mathbb{P}}
\newcommand{\E}{\mathbb{E}}
\newcommand{\eps}{\varepsilon}
\newcommand{\dd}{\,\mathrm{d}}
\begin{document}
\section*{UnsolvedMath ID: 800065}
\textbf{Source problem code:} catalogue entry
\subsection{Definitions and mathematical statement}
The supplied record does not specify the ambient hypotheses on \(M\). Its terminology suggests a smooth three-dimensional manifold, equipped with an orientation so that contact signs can be compared. Compactness, connectedness, and boundary conditions remain unspecified; the question below is conditional on the intended ambient class being clarified.

For a cooriented contact plane field \(\xi=\ker\alpha\subset TM\), the defining one-form satisfies \(\alpha\wedge\dd\alpha\ne0\). The contact structure is positive or negative according to whether this three-form agrees or disagrees with the chosen orientation of \(M\). A contact structure is tight if it contains no overtwisted disk. A homotopy of underlying plane fields is a continuous family \((\eta_t)_{t\in[0,1]}\) of rank-two subbundles of \(TM\) with prescribed endpoints; the intermediate plane fields need not be contact structures. The source's notation ``(co)oriented'' does not specify whether orientations or coorientations must be preserved by this homotopy.

A codimension-one foliation is locally a decomposition into plaques \(\R^2\times\{t\}\). It is cooriented when its transverse directions admit a consistent continuous orientation. It is Reebless when it contains no Reeb component, meaning an embedded solid torus carrying the standard Reeb foliation, whose boundary is a torus leaf and whose interior leaves are planes accumulating on that boundary. The supplied record does not specify whether \(C^0\) refers to a topological foliation or to a foliation with additional leafwise differentiability and a continuous tangent plane field.

\textbf{Question (with the qualifications above unresolved).} For every \(M\) in the intended ambient class, suppose that \((\xi_-,\xi_+)\) is a pair of (co)oriented contact structures of negative and positive sign, respectively, that both structures are tight, and that they are homotopic as plane fields in the intended sense. Must there exist a cooriented Reebless \(C^0\)-foliation \(\mathcal F\) on \(M\)?
\subsection{Short English statement}
Does the existence on \(M\) of two tight contact structures of opposite signs, homotopic as plane fields, imply the existence of a cooriented Reebless \(C^0\)-foliation? The ambient manifold hypotheses, the structure preserved by the homotopy, and the precise foliation regularity remain unspecified.
\subsection{Sources}
\begin{itemize}
\item Local UnsolvedMath registry, supplied record with permanent ID \(800065\), exact source code \texttt{catalogue entry}, in the AMR Open Problem Lists collection. The record attributes the question to Nariman and Yazdi, with proposed year 2025, and labels it open; this status is supplied metadata, not an independently verified assessment.
\item Nariman, Yazdi, \emph{Problem List on Foliations and Diffeomorphism Groups} (2025), Question 6. Supplied source URL: \url{https://arxiv.org/abs/2503.18274}. The supplied provenance records access on 2026-10-16 and extraction by \texttt{arxiv\_latex}; these provenance details are reproduced without claiming independent source access.
\item Supplied upstream revision: \texttt{d371ee917ac9647fbb06cd6fc529deec6f3dda82}. Registry creation and update timestamps are both \texttt{2026-10-16T14:00:00Z}.
\end{itemize}
\end{document}
